\documentclass[12pt,twocolumn]{article}
\usepackage[utf8]{inputenc}
\usepackage{amsmath,amssymb,amsfonts}
\usepackage{graphicx}
\usepackage{hyperref}
\usepackage{geometry}
\usepackage{booktabs}
\usepackage{siunitx}
\usepackage{float}
\usepackage{enumitem}
\usepackage{bm}
\usepackage{xcolor}

\geometry{margin=1in}

\newcommand{\grad}{\nabla}
\newcommand{\Reff}{\mathrm{eff}}
\newcommand{\mdot}{\dot{m}}
\newcommand{\vexh}{v_{\mathrm{exh}}}

\title{Atmospheric Mode of the $\psi$-Bubble Propulsion System:\\
Air as Gravitational Propellant --- Feasibility Analysis\\
and the Fundamental Density--Field Paradox}

\author{Gary Thomas Alcock\thanks{gary@gtacompanies.com}\\
\textit{Independent Researcher, Los Angeles, CA, USA}}

\date{March 2026}

\begin{document}

\twocolumn[
\maketitle

\begin{abstract}
We analyse the atmospheric operating mode of a $\psi$-bubble propulsion
device within Density Field Dynamics (DFD), where ambient air serves as
propellant accelerated gravitationally rather than thermally. The concept
is a jet engine without combustion: an asymmetric $\psi$-gradient
(Mode~II nozzle geometry) accelerates ingested air rearward via the DFD
body force $\mathbf{a} = (c^2/2)\,\kappa\,\grad\eta$. We compute mass
flow rates, exhaust velocities, and thrust for a reference 5\,m shell
with 17\,m bubble radius operating at sea level. Unconstrained estimates
yield $\sim$1.8\,GN of thrust (five Saturn~V equivalents), but realistic
energy limits reduce the exhaust velocity to $\sim$16\,km/s (Mach~47),
giving $\sim$5\,MN thrust at reduced mass flow. Critically, we identify
a \emph{fundamental density--field paradox}: the magnetic cavity that
lowers $\rho$ to reach the DFD threshold simultaneously lowers $B$ at
the air-entry boundary, so that incoming air never crosses the activation
threshold. We quantify this paradox and discuss three architectural
mitigations: nested-shell geometries, pulsed operation, and high-altitude
launch profiles.
\end{abstract}
]

%----------------------------------------------------------------------
\section{Introduction}
\label{sec:intro}
%----------------------------------------------------------------------

The key insight for atmospheric operation of a DFD $\psi$-bubble device
is that \emph{stored propellant is unnecessary}. The device uses ambient
air.

The $\psi$-bubble creates a region of modified effective refractive
index $n_\Reff$ around the device. Air molecules entering this region
experience the DFD body-force acceleration
\begin{equation}
  \mathbf{a} = \frac{c^2}{2}\,\kappa\,\grad\eta\,,
  \label{eq:body-force}
\end{equation}
where $\eta = B^2/(2\mu_0\,\rho\,c^2)$ is the dimensionless
field--density ratio and $\kappa = \alpha/4$ is the DFD vacuum loading
constant. With an asymmetric nozzle geometry (Mode~II), more air is
accelerated out the rear than is retained at the front. The device
therefore operates as a \emph{jet engine without combustion} --- the air
is accelerated gravitationally, not thermally.

Throughout this analysis we adopt a reference geometry: a
superconducting shell of radius $R_s = 5$\,m generating a surface field
of $B_0 = 20$\,T, producing a $\psi$-bubble of effective radius
$R_b = 17$\,m.

%----------------------------------------------------------------------
\section{Air Mass Flow Through the $\psi$-Bubble}
\label{sec:massflow}
%----------------------------------------------------------------------

If the device moves at velocity $V$ through the atmosphere, the mass
flow rate through the above-threshold capture area is
\begin{equation}
  \mdot = \rho_\mathrm{air}\,A_\mathrm{cap}\,V\,,
  \label{eq:mdot}
\end{equation}
where $A_\mathrm{cap} = \pi R_b^2$ is the cross-sectional area of the
above-threshold region. For the reference geometry:
\begin{equation}
  A_\mathrm{cap} = \pi \times 17^2 \approx 908~\si{m^2}\,.
\end{equation}
At sea level ($\rho_\mathrm{air} = 1.225$~\si{kg/m^3}) and
$V = 100$~\si{m/s}:
\begin{equation}
  \mdot = 1.225 \times 908 \times 100
        = 1.11 \times 10^5~\si{kg/s}\,.
  \label{eq:mdot-num}
\end{equation}
This is an enormous mass flow --- approximately 111 tonnes per second
at only 100\,m/s forward velocity.

%----------------------------------------------------------------------
\section{Exhaust Velocity of Accelerated Air}
\label{sec:vexh}
%----------------------------------------------------------------------

The air is accelerated across the $\psi$-gradient. In the idealized
(energy-unlimited) case, the maximum exhaust velocity is set by the
$\psi$-potential:
\begin{equation}
  \vexh^{(\mathrm{ideal})} = c\sqrt{\delta\psi_\Reff}\,.
  \label{eq:vexh-ideal}
\end{equation}
For $\delta\psi \sim 0.1$ (UVCS-calibrated range), this yields
$\vexh = 0.32\,c = 9.5 \times 10^7$~\si{m/s}. However, the
corresponding power requirement
\begin{equation}
  P = \tfrac{1}{2}\,\mdot\,\vexh^2
    \sim 5 \times 10^{23}~\si{W}
\end{equation}
is obviously unphysical.

\subsection{Energy-limited exhaust velocity}

The realistic exhaust velocity is set by the available magnetic energy
density $U_B = B^2/(2\mu_0)$:
\begin{equation}
  \vexh = \sqrt{\frac{2\,U_B}{\rho_\mathrm{ambient}}}\,.
  \label{eq:vexh-energy}
\end{equation}
At $B = 20$\,T:
\begin{equation}
  U_B = \frac{(20)^2}{2 \times 4\pi \times 10^{-7}}
      = 1.59 \times 10^8~\si{Pa}\,.
\end{equation}
In atmospheric air ($\rho = 1.225$~\si{kg/m^3}):
\begin{equation}
  \vexh = \sqrt{\frac{2 \times 1.59 \times 10^8}{1.225}}
        = 16{,}163~\si{m/s}
        \approx 16~\si{km/s}\,,
  \label{eq:vexh-num}
\end{equation}
corresponding to approximately Mach~47. This is physically meaningful
and represents a realistic upper bound on the air exhaust speed.

%----------------------------------------------------------------------
\section{Thrust Estimates}
\label{sec:thrust}
%----------------------------------------------------------------------

\subsection{Ideal (full-capture) thrust}

If all air through the capture area were accelerated to $\vexh$:
\begin{equation}
  F = \mdot \times \vexh
    = 1.11 \times 10^5 \times 16{,}163
    = 1.79 \times 10^9~\si{N}
    \approx 1.79~\si{GN}\,.
  \label{eq:thrust-ideal}
\end{equation}
This is 1.79 billion newtons --- the equivalent of approximately five
Saturn~V first-stage engines.

However, this assumes \emph{all} air through the capture cross-section
is accelerated to $\vexh$. In practice, only air that crosses the DFD
threshold surface is gravitationally accelerated; air remaining below
threshold simply flows past the device aerodynamically.

\subsection{Threshold-limited (cavity) thrust}
\label{sec:cavity-thrust}

Inside the magnetically evacuated cavity surrounding the shell, the
residual gas density is far below atmospheric. Taking
$\rho_\mathrm{cav} \sim 10^{-5}$~\si{kg/m^3}, the cavity mass flow at
$V = 100$\,m/s is:
\begin{equation}
  \mdot_\mathrm{cav}
    = 10^{-5} \times 908 \times 100
    = 0.91~\si{kg/s}\,.
\end{equation}
The exhaust velocity within the low-density cavity is substantially
higher:
\begin{equation}
  \vexh^{(\mathrm{cav})}
    = \sqrt{\frac{2 \times 1.59 \times 10^8}{10^{-5}}}
    = 5.66 \times 10^6~\si{m/s}
    = 0.019\,c\,.
\end{equation}
The resulting thrust is:
\begin{equation}
  F_\mathrm{cav} = 0.91 \times 5.66 \times 10^6
                  = 5.15 \times 10^6~\si{N}
                  \approx 5.15~\si{MN}\,.
  \label{eq:thrust-cavity}
\end{equation}
This is approximately 5 meganewtons --- comparable to a Space Shuttle
Main Engine (SSME). The device must continuously scoop atmospheric air
through the high-$B$ front surface, magnetically evacuate it within
the cavity, then accelerate it rearward via the threshold mechanism.

\begin{table}[t]
\centering
\caption{Summary of atmospheric thrust estimates for the reference
5\,m shell / 17\,m bubble geometry at 20\,T, $V = 100$\,m/s, sea level.}
\label{tab:thrust-summary}
\begin{tabular}{@{}lcc@{}}
\toprule
Parameter & Full capture & Cavity only \\
\midrule
$\rho$ (\si{kg/m^3}) & 1.225 & $10^{-5}$ \\
$\mdot$ (\si{kg/s}) & $1.11 \times 10^5$ & 0.91 \\
$\vexh$ (\si{m/s}) & 16{,}163 & $5.66 \times 10^6$ \\
$F$ (\si{N}) & $1.79 \times 10^9$ & $5.15 \times 10^6$ \\
Equivalent & $\sim$5 Saturn~V & $\sim$1 SSME \\
\bottomrule
\end{tabular}
\end{table}

%----------------------------------------------------------------------
\section{Threshold Analysis at Sea Level}
\label{sec:threshold}
%----------------------------------------------------------------------

The DFD threshold is crossed when $\eta/\eta_c \geq 1$, where
$\eta_c = 1.82 \times 10^{-3}$ is the critical vacuum loading ratio.
At sea level with a 20\,T field:
\begin{equation}
  \eta = \frac{B^2}{2\mu_0\,\rho_\mathrm{air}\,c^2}
       = \frac{(20)^2}{2 \times 4\pi\times10^{-7}
         \times 1.225 \times 9\times10^{16}}
       = 1.44 \times 10^{-6}\,.
\end{equation}
The threshold ratio is therefore:
\begin{equation}
  \frac{\eta}{\eta_c}
    = \frac{1.44 \times 10^{-6}}{1.82 \times 10^{-3}}
    = 7.9 \times 10^{-4}\,.
  \label{eq:eta-ratio}
\end{equation}
Sea-level air at 20\,T is \textbf{below threshold by a factor of
$\sim$1000}. The ambient atmosphere is \emph{not} above threshold.
Only the magnetically evacuated cavity, where $\rho \ll \rho_\mathrm{air}$,
can cross threshold.

%----------------------------------------------------------------------
\section{The Fundamental Density--Field Paradox}
\label{sec:paradox}
%----------------------------------------------------------------------

The cavity-thrust calculation (Section~\ref{sec:cavity-thrust})
requires the magnetic cavity to be continuously refilled with
atmospheric air at the rate gas is ejected. We now show that a
fundamental geometrical conflict prevents this in a single-shell
configuration.

\subsection{Magnetic confinement is real}

At $B = 20$\,T the magnetic pressure is:
\begin{equation}
  P_\mathrm{mag} = \frac{B^2}{2\mu_0}
                 = 1.59 \times 10^8~\si{Pa}\,,
\end{equation}
which exceeds atmospheric pressure
$P_\mathrm{atm} \approx 10^5$~\si{Pa} by three orders of magnitude.
The evacuated cavity surrounding the shell is therefore physically
robust.

\subsection{Air entry boundary}

Air can only enter the cavity where the magnetic pressure drops below
atmospheric pressure, i.e.\ where:
\begin{equation}
  B < \sqrt{2\mu_0\,P_\mathrm{atm}} \approx 0.5~\si{T}\,.
  \label{eq:B-entry}
\end{equation}
At the flight speed $V = 100$\,m/s, the ram pressure
$P_\mathrm{ram} = \frac{1}{2}\rho V^2 = 6125$\,Pa is small compared
to $P_\mathrm{mag}$, so ram pressure does not significantly aid air
entry into the high-$B$ region.

\subsection{Threshold at the entry boundary}

At the air-entry boundary ($B \sim 0.5$\,T,
$\rho \sim 1.225$~\si{kg/m^3}):
\begin{equation}
  \eta_\mathrm{entry}
    = \frac{(0.5)^2}{2 \times 4\pi\times10^{-7}
      \times 1.225 \times 9\times10^{16}}
    = 9.0 \times 10^{-10}\,.
\end{equation}
This is \emph{below threshold by a factor of $\sim$2 million}. The air
entering the cavity at the low-$B$ boundary is nowhere near threshold
activation.

Meanwhile, reaching threshold at atmospheric density would require:
\begin{equation}
  B_\mathrm{thresh}(\rho_\mathrm{air})
    = \sqrt{2\mu_0\,\rho_\mathrm{air}\,c^2\,\eta_c}
    \approx 22{,}500~\si{T}\,,
\end{equation}
which is far beyond any achievable laboratory field.

\subsection{Statement of the paradox}

\begin{quote}
\textbf{Density--Field Paradox:} The DFD threshold requires
simultaneously low $\rho$ \emph{and} high $B$. But magnetic evacuation
only produces low $\rho$ where $B$ is already low (at the cavity
boundary). By the time $B$ drops enough for air to enter, $B$ is also
too low for threshold activation. A single-shell configuration cannot
gravitationally accelerate atmospheric air.
\end{quote}

This is not a minor engineering difficulty but a fundamental
geometrical constraint of the single-shell architecture.

%----------------------------------------------------------------------
\section{Mitigation Architectures}
\label{sec:mitigations}
%----------------------------------------------------------------------

Three approaches may circumvent or weaken the density--field paradox.

\subsection{Nested multi-shell geometry}

A two-stage architecture separates the evacuation and acceleration
functions:
\begin{itemize}[nosep]
  \item \textbf{Outer shell} (moderate $B$, $\sim$1--5\,T): Creates
        a partial magnetic cavity, reducing $\rho$ by 2--3 orders of
        magnitude below atmospheric.
  \item \textbf{Inner shell} (high $B$, 20\,T): Operates in the
        pre-evacuated intermediate-density gas, where the threshold
        can be reached at achievable field strengths.
\end{itemize}
The outer shell acts as a ``pre-conditioner,'' lowering $\rho$ into a
regime where the inner shell can activate the DFD mechanism. The air
entry boundary of the outer shell is at low $B$ (as before), but
between the two shells there exists a region of both reduced $\rho$
and elevated $B$ --- precisely the conditions needed for threshold
crossing.

\subsection{Pulsed operation}

A time-domain approach exploits the difference between magnetic
evacuation timescale and gas refill timescale:
\begin{enumerate}[nosep]
  \item A high-$B$ pulse magnetically evacuates gas from the cavity
        on a fast ($\mu$s) timescale.
  \item The DFD threshold is crossed transiently before the gas
        refills on the slower (ms) acoustic/diffusion timescale.
  \item The cycle repeats at kHz rates.
\end{enumerate}
The duty cycle and time-averaged thrust depend on the ratio
$\tau_\mathrm{evac}/\tau_\mathrm{refill}$.

\subsection{High-altitude launch profile}

At altitude, ambient $\rho$ decreases exponentially while $B$ remains
fixed. The threshold field strength scales as
$B_\mathrm{thresh} \propto \sqrt{\rho}$:

\begin{table}[H]
\centering
\caption{Threshold field vs.\ altitude for DFD activation.}
\label{tab:altitude}
\begin{tabular}{@{}lcc@{}}
\toprule
Altitude & $\rho$ (\si{kg/m^3}) & $B_\mathrm{thresh}$ (\si{T}) \\
\midrule
Sea level   & 1.225          & 22{,}500 \\
10\,km      & 0.414          & 13{,}100 \\
30\,km      & 0.018          & 2{,}730 \\
50\,km      & $1.0\times10^{-3}$ & 643 \\
80\,km      & $1.8\times10^{-5}$ & 86 \\
100\,km     & $5.6\times10^{-7}$ & 15 \\
\bottomrule
\end{tabular}
\end{table}

At $\sim$100\,km altitude (K\'{a}rm\'{a}n line), a 20\,T field
\emph{exceeds} the threshold requirement. A balloon or rocket launch
to the upper mesosphere would place the device in conditions where
the single-shell architecture can operate directly in ambient gas.

This suggests a natural operational envelope: the device launches
conventionally (or via balloon) to $\sim$80--100\,km altitude, where
it activates the $\psi$-bubble atmospheric mode for acceleration
through the remaining atmosphere and into orbit.

%----------------------------------------------------------------------
\section{Discussion}
\label{sec:discussion}
%----------------------------------------------------------------------

The atmospheric mode analysis reveals both promise and a fundamental
constraint:

\paragraph{Promise.}
If the density--field paradox is resolved, the thrust levels are
extraordinary. Even the conservative cavity-only estimate of 5\,MN
matches a Space Shuttle Main Engine, and the device consumes no
chemical propellant --- only electrical power for the superconducting
magnets and any active $\psi$-field shaping.

\paragraph{Constraint.}
At sea level, no single-shell configuration at achievable field
strengths ($B \leq 45$\,T, the current laboratory record) can cross
the DFD threshold in atmospheric-density air. The paradox is
structural: the same magnetic field that evacuates the cavity also
defines the boundary conditions that prevent incoming air from
reaching threshold.

\paragraph{Path forward.}
The nested-shell and pulsed-operation architectures offer genuine
engineering paths around the paradox, but each introduces significant
complexity:
\begin{itemize}[nosep]
  \item Nested shells require independent superconducting coil systems
        with different field geometries and a structurally sound
        inter-shell region.
  \item Pulsed operation requires high-power switching at kHz rates
        for 20\,T-class coils, well beyond current pulsed-magnet
        technology for sustained operation.
  \item High-altitude launch avoids the paradox entirely but limits
        the atmospheric mode to altitudes above $\sim$80\,km.
\end{itemize}

The most practical near-term path is the high-altitude launch: a
balloon-assisted ascent to the mesosphere, followed by $\psi$-bubble
activation in the ambient low-density gas, providing enormous thrust
for orbital insertion without propellant.

%----------------------------------------------------------------------
\section{Conclusions}
\label{sec:conclusions}
%----------------------------------------------------------------------

\begin{enumerate}
  \item The $\psi$-bubble atmospheric mode concept --- using ambient
        air as gravitationally-accelerated propellant --- is physically
        sound in principle, yielding thrust levels from 5\,MN (cavity
        mode) to 1.8\,GN (full capture).

  \item The energy-limited exhaust velocity at 20\,T is
        $\sim$16\,km/s (Mach~47) in atmospheric air, or $0.019\,c$
        in the evacuated cavity.

  \item A \textbf{fundamental density--field paradox} prevents
        single-shell atmospheric operation at sea level: magnetic
        evacuation and threshold activation are geometrically
        incompatible in a single-stage architecture.

  \item Three mitigations exist: nested shells (engineering-complex),
        pulsed operation (technology-limited), and high-altitude
        launch (operationally simplest).

  \item The recommended near-term architecture is balloon launch to
        80--100\,km followed by $\psi$-bubble activation, exploiting
        the naturally low ambient density at mesospheric altitudes.
\end{enumerate}

%----------------------------------------------------------------------
% References
%----------------------------------------------------------------------
\begin{thebibliography}{9}

\bibitem{alcock2025}
G.~T. Alcock, ``Density Field Dynamics: A Complete Unified Theory,''
\textit{Preprint}, 2025.

\bibitem{alcock2025alpha}
G.~T. Alcock, ``Ab Initio Derivation of the Fine Structure Constant
from Density Field Dynamics,'' \textit{Preprint}, 2025.

\bibitem{alcock2026vacuum}
G.~T. Alcock, ``Vacuum Loading in Density Field Dynamics:
$\kappa = \alpha/4$ Derived,'' \textit{Preprint}, 2026.

\end{thebibliography}

\end{document}
